\usepackage{mathrsfs}
\let\mathcal\mathscr
\let\epsilon\varepsilon
\newcommand{\Changel}{\mathcode`l="7160}
\newcommand{\Changelback}{\mathcode`l="716C}
\newcommand{\NNoChangel}[1]{%
\expandafter\let\csname old\string#1\endcsname=#1
\let#1=\relax
\newcommand{#1}{\mathcode`l="716C\csname old\string#1\endcsname\mathcode`l="7160 }%
}

\NNoChangel{\log}
\NNoChangel{\ln}
\NNoChangel{\lim}
\NNoChangel{\limsup}
\NNoChangel{\liminf}
\NNoChangel{\varlimsup}
\NNoChangel{\varliminf}
\NNoChangel{\varinjlim}
\NNoChangel{\varprojlim}


\theoremstyle{plain}
\newtheorem{teo}{Theorem}[section]
\newtheorem{prop}[teo]{Proposition}
\newtheorem{lema}[teo]{Lemma}
\newtheorem{coro}[teo]{Corollary}

\theoremstyle{definition}
\newtheorem{obs}[teo]{Remark}
\newtheorem{ej}[teo]{Example}

\newcommand{\NN}{\mathbb{N}}
\newcommand{\CC}{\mathbb{C}}
\newcommand{\RR}{\mathbb{R}}

\newcommand{\Isom}{\mathit{Isom}}

\newcommand{\ccl}[3]{#1^{#2}(#3)}
\newcommand{\corchete}[1]{\left\{{#1}\right\}}

\DeclareMathOperator{\Char}{char}
\DeclareMathOperator{\End}{End}
\DeclareMathOperator{\mdeg}{mdeg}
\DeclareMathOperator{\Nil}{\mathcal{N}}
\DeclareMathOperator{\rank}{rank}
\DeclareMathOperator{\Tr}{Tr}
\DeclareMathOperator{\Wedge}{\mathsf{\Lambda}}

